\documentclass[11pt]{article}

\usepackage[margin=0.75cm]{geometry}
\usepackage{amsmath}
\usepackage{amssymb}
\usepackage{array}
\usepackage{enumitem}
\usepackage{svg}
\pagestyle{empty}

\setlength{\columnsep}{0.55cm}
\setlength{\parindent}{0pt}
\setlength{\parskip}{0.05cm}
\setlist[enumerate]{itemsep=0.12cm,topsep=0.08cm,leftmargin=*}
\setlist[enumerate,2]{itemsep=0.05cm,topsep=0.05cm}

% Compile-safe WaveDROM inclusion.  Before rendering, a placeholder box is shown.
\newcommand{\wavedromfig}[2][\linewidth]{%
    \IfFileExists{#2.svg}{%
        \includesvg[width=#1]{#2}%
    }{%
        \fbox{\parbox[c][2.0cm][c]{0.90\linewidth}{\centering
        Render \texttt{\detokenize{#2.json}} with\\
        \texttt{python3 render\_wavedrom.py}}}%
    }%
}

\begin{document}
\twocolumn

\begin{center}
{\Large\bfseries PHYS306/COSC330 Midterm 1}\\[0.25cm]
{\normalsize Prof. Jordan C. Hanson}\\[0.10cm]
\textbf{Name/Student ID:} \rule{5.4cm}{0.4pt}
\end{center}

\scriptsize
\textbf{Instructions:} Show enough work to make your reasoning clear. The total score is 100 points.

\section{Chapter 1: Introductory Concepts}

\begin{enumerate}
\item \textbf{[14 points]} The four traces below form a 4-bit parallel bus. Read each word as $D_3D_2D_1D_0$, with one word transferred per clock period.

\begin{center}
\wavedromfig[0.75\linewidth]{fig_1_bus}
\end{center}

The clock frequency is $f_{\mathrm{CLK}}=5.0\,\mathrm{MHz}$.

\begin{enumerate}[label=(\alph*)]
    \item Determine the clock period.
    \item How much time is represented by the eight clock periods shown?
    \item Write the eight transmitted words as hexadecimal digits.
    \item Determine the total parallel bit rate.
    \item If the same data were sent serially at the same clock frequency, what would the bit rate be?
    \item What is the duty cycle of $D_0$ over the interval shown?
\end{enumerate}
\end{enumerate}

\vspace{3.5cm}

\section{Chapter 2: Number Systems, \\ Operations, and Codes}

\begin{enumerate}
\setcounter{enumi}{1}

\item \textbf{[18 points]} Perform the following conversions and arithmetic.

\begin{enumerate}[label=(\alph*)]
    \item Convert $45_{10}$ to binary and hexadecimal.
    \item Convert $\mathrm{0x3C}$ to binary and decimal.
    \item Represent $-18_{10}$ as an 8-bit two's-complement number.
    \item Treat the following as 8-bit two's-complement numbers and add them:
    \begin{equation*}
        00101101 + 11110110.
    \end{equation*}
    Give the 8-bit result, its decimal value, and state whether signed overflow occurs.
\end{enumerate}

\vspace{5cm}

\item \textbf{[10 points]} Work with Gray code, BCD, and parity.

\begin{enumerate}[label=(\alph*)]
    \item Convert binary $1011$ to Gray code.
    \item Convert Gray code $1110$ to binary.
    \item Interpret the BCD word $0010\;0101$ as a decimal number.
    \item A transmitted word uses one leading even-parity bit followed by two BCD digits. The receiver obtains \texttt{0 0010 0101}.  Has a parity error been detected? Explain using the number of HIGH bits.
\end{enumerate}
\end{enumerate}

\vspace{3.5cm}

\section{Chapter 3: Logic Gates}

\begin{enumerate}
\setcounter{enumi}{3}

\item \textbf{[12 points]} Consider the gate network below.

\begin{center}
\wavedromfig[0.6\linewidth]{fig_4_logic}
\end{center}

\begin{enumerate}[label=(\alph*)]
    \item Write the Boolean expression for $X$.
    \item Construct the complete truth table.
    \item Simplify the expression using Boolean algebra.
    \item Which single basic logic gate performs the same function?
\end{enumerate}

\vspace{4cm}

\item \textbf{[10 points]} The data below are inputs to the gate shown.

\begin{center}
\wavedromfig[0.6\linewidth]{fig_5_inputs}\\[0.10cm]
\wavedromfig[0.5\linewidth]{fig_5_gate}
\end{center}

\begin{enumerate}[label=(\alph*)]
    \item Draw the output waveform for $X$.
    \item If the bit interval is $20\,\mathrm{ns}$, give the width of the LOW pulse in $X$.
\end{enumerate}

\vspace{4cm}

\item \textbf{[12 points]} NAND gates are universal logic elements.

\begin{enumerate}[label=(\alph*)]
    \item Using only 2-input NAND gates, draw a circuit that implements a 2-input OR gate.
    \item We apply $A=1$ and $B=0$ to a 2-input NAND gate, and the output is $X=0$.  Is that gate behaving correctly?  Why or why not?
\end{enumerate}
\end{enumerate}

\vspace{4cm}

\section{Chapter 4: Boolean Algebra and Logic Simplification}

\begin{enumerate}
\setcounter{enumi}{6}

\item \textbf{[14 points]} Simplify the following logic functions.

\begin{enumerate}[label=(\alph*)]
    \item Apply DeMorgan's theorems to
    \begin{equation*}
        Y=\overline{(A+\overline{B})(C+\overline{D})}.
    \end{equation*}

    \item Use Boolean algebra to reduce
    \begin{equation*}
        Z=A\overline{B}+AB+\overline{A}B
    \end{equation*}
    to its minimum form.

    \item Use a 4-variable Karnaugh map to find the minimum SOP form of
    \begin{align*}
    X={}&\overline{A}\,\overline{B}\,\overline{C}D
      +\overline{A}\,\overline{B}CD
      +\overline{A}B\overline{C}D
      +\overline{A}BCD\\
      &+AB\overline{C}\,\overline{D}
      +AB\overline{C}D
      +ABC\overline{D}
      +ABCD.
    \end{align*}
    Draw the minimum circuit using basic gates.
\end{enumerate}
\end{enumerate}

\vspace{6cm}

\section{Chapter 5: Combinational Logic \\ Analysis}

\begin{enumerate}
\setcounter{enumi}{7}

\item \textbf{[10 points]} A receiver accepts three data bits $A,B,C$ and a received parity bit $P$. The transmitter uses \textbf{odd parity}. Define an error output $E$ so that $E=0$ when the received four-bit word has correct odd parity, and $E=1$ when the parity is incorrect.

\begin{enumerate}[label=(\alph*)]
    \item Write a Boolean expression for $E$ using XOR/XNOR notation.
    \item For $ABC=101$, determine the parity bit $P$ that should be transmitted.
    \item If $ABC=101$ is received with $P=0$, determine $E$.
    \item Draw the minimum gate-level circuit for the parity checker.
\end{enumerate}
\end{enumerate}

\end{document}
